\documentclass[11pt]{article}
\usepackage[utf8]{inputenc}
\usepackage[margin=1in]{geometry}
\usepackage{amsmath,amssymb}
\usepackage{booktabs}
\usepackage{graphicx}
\usepackage{xcolor}
\usepackage[colorlinks=true,linkcolor=blue,citecolor=blue,urlcolor=blue]{hyperref}
\usepackage{authblk}

% ============================================================================
% WORKSHOP PAPER 2 -- quantization-technique angle
% Real measurements: NeuroNCAP closed-loop (real-camera), scenes 0103 frontal + 0796 stationary, 10 seeds each.
% Primary metric = generation coherence (well-formed predicted-frame rate).
% ============================================================================

\title{\textbf{Granularity, Not Bit-Width: Group-Wise INT4 Preserves Generation\\
Coherence of Driving-VLA Planners Under Deployment-Domain Shift}}
\author[1]{Justin Williams}
\author[1]{Kishor Datta Gupta}
\author[1]{Roy George}
\author[1]{Mrinmoy Sarkar}
\affil[1]{Clark Atlanta University \quad \texttt{justinwilliamstech693@gmail.com}}
\date{}

\begin{document}
\maketitle

\begin{abstract}
We study weight-only INT4 quantization of the language/action head of OpenDriveVLA-0.5B in a
\emph{closed-loop} neural-rendering simulator (NeuroNCAP, two real-camera scenes). The headline: under the
deployment-domain shift induced by photorealistic re-rendering, \emph{naive per-output-channel}
INT4 catastrophically loses generation coherence---the planner emits degenerate trajectory text on
nearly every frame and effectively freezes ($0.0\%$ and $8.6\%$ well-formed across two scenes)---
whereas \emph{group-wise} INT4 (group size 128, an extra $\approx 0.1$ bits/weight) recovers most of
FP16's coherence ($72.9\%$ vs $74.3\%$; $57.9\%$ vs $72.4\%$). The same per-channel INT4 is lossless
and well-formed \emph{in-distribution} (open-loop, real images), so the failure is a
quantization\,$\times$\,domain-shift interaction governed by \emph{scale granularity}, not by 4-bit
precision per se. We report generation-coherence rate as the clean metric and discuss why downstream
collision/NCAP counts are a confounded, scene-dependent proxy (on one scene every config including
FP16 collides).
\end{abstract}

\section{Introduction}
Group-wise (block-wise) weight quantization with per-group scales is the de-facto standard for
INT4 LLMs because a single scale shared across many weights is set by the group's largest
magnitude, so one outlier inflates the quantization step for every weight it shares it with. We
ask whether this granularity matters for the \emph{closed-loop} behavior of a quantized driving
VLA under the distribution shift of a neural-rendering deployment surrogate, and find that it is
decisive.

\section{Setup}
We quantize only the 168 Linear layers of the Qwen2.5-0.5B backbone (357.8\,M parameters);
perception (UniAD track+map), projectors, embeddings, and head stay FP. We run closed-loop on
two NeuroNCAP scenes (0103 frontal, 0796 stationary): a NeuRAD renderer re-renders six cameras at
the ego's \emph{driven} pose each step, a model node returns a trajectory, and a kinematic controller
advances the vehicle. Ten random actor-jitter seeds are run \emph{paired} across precisions
(jitter seeded by run index; only the quantization differs). Our primary metric is
\textbf{generation coherence}: the fraction of frames on which the planner emits a well-formed
(non-degenerate) trajectory, computed from the logged per-frame plans.

\section{Results}
\begin{table}[h]\centering
\caption{Closed-loop generation coherence (well-formed predicted-frame rate), \emph{real-camera}
NeuRAD rendering, 10 paired seeds per scene. Weight-only quantization of the backbone.}
\begin{tabular}{lccc}
\toprule
Config & 0103 frontal & 0796 stationary & recovers FP16? \\
\midrule
FP16                       & 74.3\%          & 72.4\%          & --- \\
INT4 per-channel (naive)   & \textbf{0.0\%}  & \textbf{8.6\%}  & no (total collapse) \\
INT4 group-128             & \textbf{72.9\%} & \textbf{57.9\%} & yes / partial \\
\bottomrule
\end{tabular}
\end{table}

\paragraph{The effect is domain-shift--induced, not 4-bit precision.} On real nuScenes images
(open-loop) the \emph{same} per-channel INT4 is lossless (0.33\,m L2, well-formed). It collapses
only on rendered (out-of-distribution) inputs, and this replicates across both scenes. Group-wise
INT4 is robust on both (fully recovering FP16 on 0103, partially on 0796). The fix---finer scale
granularity---costs $16/128 = 0.125$ extra bits per weight.

\paragraph{Mechanism.} A shared INT4 scale is $s=\max|w|/7$. One large-magnitude weight inflates
$s$ for the whole group, coarsening the grid for the small weights that carry signal. Group-wise
scales confine each outlier to its own block. Under distribution shift the planner's hidden states
move; per-channel's coarsened grid then tips greedy decoding into malformed trajectory text, while
group-wise's finer grid preserves coherence.

\section{On the metric}
We deliberately report generation coherence rather than collision rate. The downstream collision
outcome is a \emph{confounded, scene-dependent proxy}: on 0103 per-channel INT4 collides on 4/10
seeds vs group-wise/FP16 on 0/10, but on 0796 \emph{every} config---including FP16---collides 10/10
because the base policy cannot clear that scenario. A collapsed (frozen) planner only collides on
geometries that strike a stationary ego, and the count depends on the degenerate fallback. Coherence
isolates the quantization effect; collisions do not.

\section{Limitations}
Results are on two scenes; the rendered inputs are out-of-distribution to real nuScenes (even FP16
fails scene-0796's collision test, $10/10$), so the claim is the relative one (group $\approx$ FP16
$\gg$ per-channel on coherence). Group-wise recovery is \emph{partial} on the harder scene ($57.9\%$
vs FP16 $72.4\%$)---we do not claim exact equivalence. Lower-bit (INT3/INT2) and activation-aware
(AWQ) variants were characterized only in an earlier single-scene setup and are not re-reported here
pending a real-camera re-run.

\section{Conclusion}
For quantized driving-VLA planners under deployment-domain shift, \emph{scale granularity}, not
bit-width, governs whether generation stays coherent. Group-wise INT4 ($\approx 0.1$ extra
bits/weight) matches FP16 robustness where naive per-channel INT4 collapses---an actionable,
nearly-free deployment recommendation that open-loop L2 cannot motivate.

\begin{thebibliography}{9}
\bibitem{awq} Lin et al. \emph{AWQ: Activation-aware Weight Quantization for LLM Compression}. MLSys 2024.
\bibitem{gptq} Frantar et al. \emph{GPTQ: Accurate Post-Training Quantization for GPTs}. ICLR 2023.
\bibitem{neuroncap} NeuroNCAP. ECCV 2024.
\bibitem{opendrivevla} OpenDriveVLA. 2025. (Replace with verified citation.)
\end{thebibliography}
\end{document}
